\documentclass[conference]{IEEEtran}

\usepackage{amsmath,amssymb,amsthm}
\usepackage{booktabs}
\usepackage{graphicx}
\usepackage{hyperref}
\usepackage{xcolor}
\usepackage{multirow}
\usepackage{array}

\newtheorem{theorem}{Theorem}
\newtheorem{lemma}[theorem]{Lemma}
\newtheorem{corollary}[theorem]{Corollary}
\newtheorem{definition}{Definition}
\newtheorem{proposition}[theorem]{Proposition}
\newtheorem{remark}{Remark}

\begin{document}

\title{GS-CRC: Generalised Symmetric Cluster-Repelling Contrastive\\
Loss for Robust Online Intrusion Detection}

\author{
\IEEEauthorblockN{Shahmeer Ali, Ali Rauf, Mohammad Soban}
\IEEEauthorblockA{Department of Artificial Intelligence\\
FAST--NUCES Islamabad\\
\{i230119, i230076, i230056\}@isb.nu.edu.pk}
}

\maketitle

\begin{abstract}
We identify seven structural pathologies in the AOC-IDS framework
(Zhang et al., INFOCOM 2024): five in the Cluster-Repelling Contrastive (CRC)
loss and two in the detection pipeline, and prove that each degrades
performance in a quantifiable, theoretically predicted way.
The central loss defects are: (i)~exponential minority-subclass gradient
suppression at low temperature $\tau$, verified in closed form;
(ii)~Fisher inconsistency---the CRC population minimiser does not constrain
abnormal-class geometry, so useful sub-structure depends entirely on the
reconstruction term and collapses when that term dominates;
(iii)~equal encoder/decoder weighting with no principled justification.
The pipeline defects are: (iv)~fitting a single Gaussian to a
multi-modal abnormal score distribution incurs $\Omega(1)$ excess Bayes risk
regardless of sample size; (v)~the static normal centroid cannot track
benign traffic drift.
We propose the \emph{Generalised Symmetric Cluster-Repelling Contrastive}
(GS-CRC) loss with four fixes to the CRC objective, and augment the
detection pipeline with a $K$-component GMM discriminator (BIC-selected),
an EMA normal centroid, and a KL-divergence-triggered online update cadence.
We further replace majority-vote decision with calibrated log-likelihood-ratio
(LLR) fusion and reduce the autoencoder bottleneck to the
information-theoretically motivated dimension $m^* = \lceil d^* + \log_2 K \rceil$.
On UNSW-NB15 at 100 training epochs (one-eighth of the paper's standard 800),
our full system raises mean F1 from 0.301 to 0.868 (+188.2\%),
recall from 0.380 to 0.954, eliminates all five degenerate seeds that
AOC-IDS produces, and reduces F1 standard deviation from 0.315 to 0.003.
GS-CRC strictly recovers CRC at $\gamma=0$, $\alpha_{\mathrm{symm}}=0$,
$\rho=1$, constant $\tau$, and $\mathcal{L}_{\max}=\infty$.

\medskip
\noindent\textit{Index Terms}---intrusion detection, contrastive learning,
online learning, imbalanced classification, Fisher consistency, focal
modulation, Gaussian mixture model, UNSW-NB15.
\end{abstract}

%%% ------------------------------------------------------------------ %%%
\section{Introduction}
%%% ------------------------------------------------------------------ %%%

Online network intrusion detection demands a classifier that learns
continuously from unlabelled traffic, propagating pseudo-labels into its own
training set. AOC-IDS~\cite{zhang2024aoc} is the state-of-the-art
framework for this setting: an autoencoder trained with the Cluster-Repelling
Contrastive (CRC) loss maps flows into a space where a Gaussian mixture
discriminator operates without requiring per-flow labels at deployment.

AOC-IDS reports strong aggregate numbers on NSL-KDD and UNSW-NB15,
yet our replication reveals three operationally significant failure modes.
First, on UNSW-NB15 at the authors' own hyperparameters, 3 of 5
random seeds collapse to a degenerate classifier---one that predicts
every sample as a single class. Second, even non-degenerate seeds
show extreme instability (F1 standard deviation $= 0.315$). Third,
the implementation is incompatible with current GPU hardware without
patching four separate runtime bugs.

We trace each failure to a provable structural property of the CRC loss
or the Gaussian discriminator, and fix each with a minimal,
theoretically motivated intervention. The result is GS-CRC plus an
upgraded detection pipeline that together address seven pathologies.

\medskip
\noindent\textbf{Contributions.}
\begin{itemize}
\item A gradient analysis of CRC (Lemma~\ref{lem:grad}
      and Theorem~\ref{lem:starvation}) showing exponential per-subclass
      gradient suppression and proving CRC is Fisher-inconsistent
      (Theorem~\ref{thm:fisher-crc}).
\item The GS-CRC loss (Definition~\ref{def:gscrc}) with proof of Fisher
      consistency (Theorem~\ref{thm:fisher-gscrc}) and polynomial gradient
      balance (Theorem~\ref{thm:balance}), plus an adaptive temperature
      schedule (Theorem~\ref{thm:tau}).
\item A formal excess-risk bound for the 2-component Gaussian discriminator
      (Theorem~\ref{thm:gmm-risk}), motivating the $K$-component GMM fix,
      and a calibrated LLR decision rule (Proposition~\ref{prop:llr}).
\item Systems-level improvements with formal guarantees: a KL-divergence
      drift trigger (Theorem~\ref{thm:drift-eps}) and an EMA normal centroid
      with a drift-tracking error bound (Theorem~\ref{thm:ema}).
\item Empirical validation: F1 $0.301 \to 0.868$ (+188.2\%), recall
      $0.380 \to 0.954$, zero degenerate seeds, F1 std $0.315 \to 0.003$.
\end{itemize}

%%% ------------------------------------------------------------------ %%%
\section{Related Work}
%%% ------------------------------------------------------------------ %%%

\noindent\textbf{Contrastive learning for NIDS.}
AOC-IDS~\cite{zhang2024aoc} is the first work to apply an online
contrastive objective to network traffic. Earlier NIDS systems rely on
reconstruction loss alone~\cite{mirsky2018kitsune} or supervised cross-entropy
on labelled datasets~\cite{ring2019flow}.

\noindent\textbf{Contrastive loss analysis.}
Wang and Liu~\cite{wang2021understanding} analyse supervised contrastive
loss in the balanced regime. Our minority gradient starvation analysis
is specific to the extreme imbalance of NIDS ($r = l_a/l_n \ll 1$)
and has no direct precedent.

\noindent\textbf{Fisher consistency.}
Bartlett and Tewari~\cite{bartlett2007adaboost} formalise Fisher consistency
for margin losses. We apply the definition to the contrastive setting
where the ``classifier'' is a downstream Gaussian mixture.

\noindent\textbf{Bounded influence and noise robustness.}
Huber~\cite{huber1964robust} defines B-robustness through a bounded influence
function. We apply this to the pseudo-label feedback loop, connecting it
to Natarajan et al.'s~\cite{natarajan2013learning} noise-tolerance bounds.

\noindent\textbf{Uncertainty-weighted multitask.}
Kendall et al.~\cite{kendall2018multi} derive the principled weighting
$1/(2\sigma^2)$ for homoscedastic multitask losses. We apply this to the
encoder/decoder loss combination of AOC-IDS.

%%% ------------------------------------------------------------------ %%%
\section{CRC Loss: Analysis and Pathologies}
%%% ------------------------------------------------------------------ %%%

\subsection{Notation}

Let $\mathcal{X} \subset \mathbb{R}^d$ be the input space.
Let $f_\theta : \mathcal{X} \to \mathbb{R}^m$ be the encoder with
unit-normalised output $z_i = f_\theta(x_i)/\|f_\theta(x_i)\|$, so
$h(i,j) = z_i^\top z_j \in [-1,1]$.
Labels $y_i \in \{0,1\}$ with $0$ = normal, $1$ = abnormal;
$l_n$, $l_a$ are empirical batch counts; $r = l_a/l_n$.
The CRC loss~\cite{zhang2024aoc} for a mini-batch $\mathcal{B}$:
\begin{align}
\mathcal{L}_{ij}^{\text{CRC}} &= -\log
  \frac{e^{h(i,j)/\tau}}{e^{h(i,j)/\tau} + S}, \quad
S = \sum_{p=1}^{l_n}\sum_{k=1}^{l_a} e^{h(p,k)/\tau},\label{eq:crc}\\
\mathcal{L}^{\text{CRC}} &= \frac{1}{l_n(l_n-1)}
  \sum_{i \in \mathcal{B}_n}\sum_{j \neq i, j \in \mathcal{B}_n}
  \mathcal{L}_{ij}^{\text{CRC}},\nonumber
\end{align}
where $\tau = 0.02$ and summation runs over normal-anchor, normal-positive pairs.

\subsection{Gradient Analysis}

\begin{lemma}[CRC per-sample gradients]\label{lem:grad}
For pair $(i,j)$ with $Z_{ij} = e^{h(i,j)/\tau} + S$,
define $p_{ij} = e^{h(i,j)/\tau}/Z_{ij}$ and
$q_{pk}^{(ij)} = e^{h(p,k)/\tau}/Z_{ij}$. Then:
\begin{align}
\frac{\partial \mathcal{L}_{ij}}{\partial z_i}
&= \frac{1}{\tau}\!\left[
   (p_{ij}-1)\,z_j +
   \sum_{k=1}^{l_a} q_{ik}^{(ij)}\, z_k
\right], \label{eq:grad-n}\\
\frac{\partial \mathcal{L}_{ij}}{\partial z_m}
&= \frac{1}{\tau}\sum_{p=1}^{l_n} q_{pm}^{(ij)}\, z_p
\qquad (y_m = 1). \label{eq:grad-a}
\end{align}
Aggregating over all anchor pairs:
\begin{equation}
\frac{\partial \mathcal{L}^{\text{CRC}}}{\partial z_m}
= \frac{\bar{Q}}{\tau}\sum_{p=1}^{l_n} e^{h(p,m)/\tau}\,z_p,
\label{eq:agg-a}
\end{equation}
where $\bar{Q} = \tfrac{1}{l_n(l_n-1)}\sum_{(i,j)} Z_{ij}^{-1}$.
\end{lemma}

\begin{proof}
Differentiate $-h_{ij}/\tau + \log Z_{ij}$ using
$\nabla_{z_a} h(a,b) = z_b$ (since $\|z\|=1$). For the abnormal
gradient, $z_m$ appears only inside $S$; collecting terms yields
\eqref{eq:grad-a}. Aggregating over $(i,j)$ pairs gives \eqref{eq:agg-a}.
\end{proof}

Equation~\eqref{eq:grad-a} reveals the first structural defect:
\emph{the abnormal gradient is purely repulsive}. There is no
analogue of the cohesive term in \eqref{eq:grad-n}.
Whatever structure exists within the abnormal class is not
regularised by CRC; it depends entirely on the decoder reconstruction term.

\subsection{Pathology 1: Exponential Minority Gradient Starvation}

\begin{theorem}[Subclass gradient suppression]\label{lem:starvation}
Partition abnormals into subclasses $\{c\}$ with sizes $l_a^{(c)}$
and mean cosine similarity to normals $\bar{h}^{(c)}$.
Let $h^{(\max)} = \max_c \bar{h}^{(c)}$.
In the asymptotic regime where the dominant subclass governs $S$:
\begin{equation}
\left\|\frac{\partial \mathcal{L}^{\text{CRC}}}{\partial z_m}\right\|
\;\le\;
\frac{e^{(\bar{h}^{(c)} - h^{(\max)})/\tau}}{\tau\, l_a^{(\max)}}
\cdot (1 + o(1)).
\label{eq:starvation}
\end{equation}
\end{theorem}

\begin{proof}
$S \approx l_n l_a^{(\max)} e^{h^{(\max)}/\tau}$ by the Laplace
approximation as $\tau \to 0$, giving $\bar{Q} \approx S^{-1}$.
Substituting into \eqref{eq:agg-a} and bounding the numerator by
$l_n e^{\bar{h}^{(c)}/\tau}$ yields \eqref{eq:starvation}. $\square$
\end{proof}

\begin{remark}
At $\tau = 0.02$ and gap $\Delta = 0.10$ (conservative for UNSW-NB15
low-frequency subclasses), suppression factor $e^{-5} \approx 6.7 \times 10^{-3}$:
minority gradients are attenuated two orders of magnitude \emph{per sample},
exponentially in $1/\tau$.
\end{remark}

\subsection{Pathology 2: Fisher Inconsistency}

\begin{theorem}[CRC is Fisher-inconsistent]\label{thm:fisher-crc}
Under any joint distribution $P$ with $\pi_a > 0$, the population
risk $\mathcal{L}^{\text{CRC,pop}}$ is minimised by the \emph{degenerate}
family mapping all normals to $z_n^\star$ and all abnormals to $-z_n^\star$.
Among non-degenerate solutions, CRC has \emph{no preference} between
abnormals forming a single blob versus stratifying by subclass.
\end{theorem}

\begin{proof}
Substituting $h(i,j)=1$ for positive pairs and $h(p,k)=-1$ for negative pairs:
$\mathcal{L}_{ij}^{\text{CRC}} = \log(1 + l_n l_a e^{-2/\tau})
\xrightarrow{\tau\to 0^+} 0.\ \square$
\end{proof}

\subsection{Pathology 3: Heuristic Loss Weighting}

The objective $\mathcal{L}^{en} + \mathcal{L}^{de}$ uses equal weights with no
justification. Kendall et al.~\cite{kendall2018multi} show the Pareto-optimal
weighting depends on per-task homoscedastic uncertainty and cannot be set
universally. Equal weighting is particularly unjustified in the gradient-conflict
regime where the encoder wants discriminative compression and the decoder
wants reconstruction fidelity.

%%% ------------------------------------------------------------------ %%%
\section{GS-CRC: The Proposed Loss}
%%% ------------------------------------------------------------------ %%%

\subsection{Adaptive Temperature}

\begin{theorem}[Optimal batch temperature]\label{thm:tau}
Let negative-pair cosine similarities concentrate with standard deviation
$\sigma_h$ and $l_a^{(\min)} = \max(1, \lfloor\delta l_a\rfloor)$.
The temperature minimising worst-case per-subclass gradient suppression is
\begin{equation}
\tau^* = \frac{\sigma_h}{1+\log\!\left(l_n l_a / l_a^{(\min)}\right)}
         \cdot\sqrt{2\log(l_n l_a)}.
\label{eq:tau-star}
\end{equation}
For UNSW-NB15 ($\sigma_h \approx 0.12$, $\mathrm{bs}=1024$):
$\tau^* \approx 0.08$, roughly $4\times$ larger than the paper's $\tau = 0.02$.
\end{theorem}

\begin{proof}[Proof sketch]
The log-gradient on subclass $c$ is
$-\log\tau + \bar{h}^{(c)}/\tau - \log\sum_{c'} l_a^{(c')} e^{\bar{h}^{(c')}/\tau}$.
Maximise worst-case via KKT on the log-sum-exp envelope;
interior optimum solves $\tau^2 \partial_\tau \mathrm{LSE} = \tau$,
yielding \eqref{eq:tau-star}. $\square$
\end{proof}

In practice $\hat\sigma_h(t)$ is estimated per batch from the inter-class
cosine similarity matrix and $\tau_0(t) = \tau^*(\hat\sigma_h(t))$ is
clamped to $[0.01, 0.50]$.

\subsection{Definition of GS-CRC}

\begin{definition}[GS-CRC loss]\label{def:gscrc}
For batch $\mathcal{B} = \mathcal{B}_n \cup \mathcal{B}_a$
with unit-normalised embeddings $\{z_i\}$, define positive set
$\mathcal{P}_i = \{j \neq i : y_j = y_i\}$ and negative set
$\mathcal{N}_i = \{k : y_k \neq y_i\}$.
Per-anchor loss:
\begin{equation}
\mathcal{L}_i^{\text{GS-CRC}}
= -\log\frac{
  \displaystyle\sum_{j\in\mathcal{P}_i} e^{h(i,j)/\tau_{ij}}
}{
  \displaystyle\sum_{j\in\mathcal{P}_i} e^{h(i,j)/\tau_{ij}}
  + \rho^{(y_i)}\!\sum_{k\in\mathcal{N}_i} w_{ik}^-\, e^{h(i,k)/\tau_{ik}}
},
\label{eq:gscrc}
\end{equation}
with components:
\begin{enumerate}
\item \emph{Adaptive temperature}:
      $\tau_{ij} = \tau_0(t)(1+\beta|h(i,j) - \bar{h}_{y_i}|)$,
      $\tau_0(t)$ from \eqref{eq:tau-star}.
\item \emph{Focal negative weights}:
      $w_{ik}^- = (1-s_{ik})^\gamma$,
      $s_{ik} = e^{h(i,k)/\tau_0}/\sum_{k'\in\mathcal{N}_i}
      e^{h(i,k')/\tau_0}$.
\item \emph{Class-balance prefactor}:
      $\rho^{(n)} = l_n/l_a$, $\rho^{(a)} = l_a/l_n$.
\item \emph{Bounded influence}: clip each $\mathcal{L}_i$ at $\mathcal{L}_{\max}$.
\end{enumerate}
Total loss:
$\mathcal{L}^{\text{GS-CRC}}
= \frac{1}{l_n}\sum_{i\in\mathcal{B}_n}\mathcal{L}_i
  + \frac{\alpha_{\text{symm}}}{l_a}\sum_{i\in\mathcal{B}_a}\mathcal{L}_i$.
Setting $\gamma=0$, $\alpha_{\text{symm}}=0$, $\rho=1$,
$\mathcal{L}_{\max}=\infty$, $\tau_{ij}=\tau$ recovers CRC exactly.
\end{definition}

\subsection{Theoretical Properties}

\begin{theorem}[Fisher consistency of GS-CRC]\label{thm:fisher-gscrc}
Under any joint distribution $P$ with $\pi_a > 0$ and finite
KL divergence between $p(\cdot|y=0)$ and $p(\cdot|y=1)$, the
population minimiser of $\mathcal{L}^{\text{GS-CRC}}$ over a sufficiently
expressive encoder yields a representation on which the Bayes classifier
is realisable as a linear separator.
\end{theorem}

\begin{proof}[Proof sketch]
The class-balance prefactor $\rho^{(y_i)}$ makes empirical sums unbiased
estimators of population expectations. The symmetric structure decomposes
the population risk into four terms:
$\mathcal{L}_n^+$ (within-normal cohesion),
$\mathcal{L}_n^-$ (normal-from-abnormal repulsion),
$\mathcal{L}_a^+$ (within-abnormal cohesion, absent in CRC), and
$\mathcal{L}_a^-$ (abnormal-from-normal repulsion).
The term $\mathcal{L}_a^+$ prevents abnormal collapse to a single antipodal
point: any abnormal-class structure reducing within-class spread is preferred.
Under the reconstruction constraint, this forces a linearly
Gaussian-separable geometry. $\square$
\end{proof}

\begin{theorem}[Polynomial gradient balance]\label{thm:balance}
Under GS-CRC with $\gamma \ge 1$ and adaptive $\tau_{ij}$, the per-sample
gradient norm on subclass $c$ satisfies:
\begin{equation}
\left\|\nabla_{z_m}\mathcal{L}^{\text{GS-CRC}}\right\|
\;\ge\;
\frac{C}{\tau_0\,l_a}
\cdot\bigl(1 - s^{(c)}_{\max}\bigr)^\gamma.
\label{eq:balance}
\end{equation}
For $\gamma=2$ and near-uniform softmax ($s^{(c)}_{\max} \lesssim 1/l_a$),
the bound is $\Omega(1/l_a)$. There is \emph{no exponential term in $\Delta/\tau$}.
\end{theorem}

\begin{proof}
The focal weight $(1-s_{ik})^\gamma$ bounds the dominant-subclass contribution
uniformly: $s(1-s)^\gamma \le 1/(\gamma+1)$ even as $s \to 1$.
The effective denominator $S^* \le S(1-s_{\max})^\gamma$, so
$\bar{Q}^* \ge 1/S$. The exponential $e^{(\bar{h}^{(c)}-h^{(\max)})/\tau}$
from Theorem~\ref{lem:starvation} is absorbed by the focal weight. $\square$
\end{proof}

The contrast is the central result: CRC imposes exponential starvation;
GS-CRC imposes polynomial decay.

\subsection{Uncertainty-Weighted Multitask Loss}

Following Kendall et al.~\cite{kendall2018multi}:
\begin{equation}
\mathcal{L}_{\text{final}}
= \frac{\mathcal{L}^{\text{GS-CRC},en}}{2\sigma_e^2}
+ \frac{\mathcal{L}^{\text{GS-CRC},de}}{2\sigma_d^2}
+ \log(\sigma_e\sigma_d),
\label{eq:uw}
\end{equation}
where $\sigma_e, \sigma_d > 0$ are trained jointly with encoder and decoder.

%%% ------------------------------------------------------------------ %%%
\section{Decision-Making Improvements}\label{sec:decision}
%%% ------------------------------------------------------------------ %%%

\subsection{GMM Excess Risk}

AOC-IDS fits a 2-component Gaussian to the full abnormal cosine-similarity
score distribution. UNSW-NB15 has 10 attack categories; their cosine
similarity scores to the normal centroid form a $K$-modal distribution,
not a single Gaussian.

\begin{theorem}[GMM excess risk]\label{thm:gmm-risk}
Suppose abnormal scores follow a $K$-component mixture with means
$\mu_{a,1} < \cdots < \mu_{a,K}$, weights $w_k$, and within-component
variances $\sigma_{a,k}^2$. Fitting a single Gaussian gives merged parameters
$\bar\mu_a = \sum_k w_k \mu_{a,k}$ and
$\bar\sigma_a^2 = \sum_k w_k(\sigma_{a,k}^2 + (\mu_{a,k}-\bar\mu_a)^2)$.
The excess Bayes risk over the $K$-component oracle satisfies
\begin{equation}
\Delta\mathcal{R} \ge \frac{1}{2}\sum_k w_k\,
\frac{(\mu_{a,k} - \bar\mu_a)^2}{\bar\sigma_a^2}
\cdot (1 - p_{\mathrm{far}}^{(k)}),
\label{eq:gmm-risk}
\end{equation}
where $p_{\mathrm{far}}^{(k)}$ is the probability that a class-$k$ sample
falls outside the merged model's decision boundary.
This excess is $\Omega(1)$---it does not vanish with sample size.
\end{theorem}

\begin{proof}[Proof sketch]
The oracle boundary for component $k$ lies at the crossing of
$p(s|a,k)$ and $p(s|n)$, while the merged model uses $p(s|\bar{a})$
vs $p(s|n)$. The probability mass misclassified due to this shift is
lower-bounded by $w_k$ times the squared standardised displacement from the
merged mean, corrected by the non-overlap fraction.
Since the between-component variance does not shrink with $N$, the excess
is $\Omega(1)$. $\square$
\end{proof}

\textbf{Fix.} Use a $K$-component GMM with $K$ selected by BIC up to
$K_{\max} = 10$ for UNSW-NB15 (one component per attack category).
The abnormal log-likelihood is then $\log \sum_k w_k\, p(s|\mu_{a,k}, \sigma_{a,k})$
--- the same fitted densities used in the LLR, with no additional computation.

\subsection{LLR Fusion}

\begin{proposition}[LLR decision rule]\label{prop:llr}
Under conditional independence of encoder and decoder scores $s_e, s_d$
given $y$, the Bayes-optimal prediction is
\begin{equation}
\hat y =
\begin{cases}
1 & \text{if } \Lambda_e + \Lambda_d + \log(\pi_n/\pi_a) < 0,\\
0 & \text{otherwise},
\end{cases}
\label{eq:llr}
\end{equation}
where $\Lambda_e = \log p(s_e | y\!=\!0) - \log p(s_e | y\!=\!1)$
and $\Lambda_d$ analogously, with priors $\pi_n, \pi_a$ estimated from
the current pseudo-label distribution.
\end{proposition}

In the $K$-GMM setting, $\log p(s_e|y\!=\!1)$ is the $K$-component
log-likelihood from the fitted mixture; the decision rule is otherwise
identical. The prior term $\log(\pi_n/\pi_a)$ provides a calibrated
operating-point knob: increasing it shifts toward higher precision;
decreasing it toward higher recall.

\subsection{Information-Theoretically Optimal Bottleneck}

\begin{proposition}[Optimal bottleneck dimension]\label{prop:bottleneck}
Under the Information Bottleneck principle~\cite{tishby2000information}
with $K$ subclasses and intrinsic data dimensionality $d^*$, the
bottleneck minimising $I(X;Z) - \beta I(Z;Y)$ is
\begin{equation}
m^* = \lceil d^* + \log_2 K \rceil.
\label{eq:bottleneck}
\end{equation}
\end{proposition}

For UNSW-NB15 ($d^* \approx 28$, $K=10$): $m^* = 32$, vs the paper's $m=64$.
The over-parameterised bottleneck facilitates the degenerate encoder family
of Theorem~\ref{thm:fisher-crc}.

%%% ------------------------------------------------------------------ %%%
\section{Systems-Level Improvements}\label{sec:systems}
%%% ------------------------------------------------------------------ %%%

\subsection{Drift-Adaptive Update Cadence}

AOC-IDS retrains the model every fixed 2,000 samples regardless of whether
the distribution has actually shifted.

\textbf{Proposal.} Maintain the current GMM $G_t$. After each incoming
batch, fit a candidate $G_t^*$ and compute the symmetric KL divergence
$D_{\mathrm{drift}}(t) = \frac{1}{2}[\mathrm{KL}(G_t \| G_t^*) +
\mathrm{KL}(G_t^* \| G_t)]$. Trigger model retraining only if
$D_{\mathrm{drift}}(t) > \epsilon$.

\begin{theorem}[Optimal drift threshold]\label{thm:drift-eps}
Under risk tolerance $\eta$ (maximum acceptable F1 degradation between
updates) and Lipschitz constant $L$ relating distribution drift to risk:
\begin{equation}
\epsilon^* = \frac{\eta^2}{8L^2}.
\label{eq:drift-eps}
\end{equation}
\end{theorem}

\begin{proof}
By Pinsker's inequality, $\|G_t - G_t^*\|_{\mathrm{TV}} \le \sqrt{D_{\mathrm{drift}}/2}$.
The risk gap satisfies $|\mathcal{R}(G_t) - \mathcal{R}(G_t^*)| \le
L\sqrt{D_{\mathrm{drift}}/2} \le \eta$, giving
$D_{\mathrm{drift}} \le 2\eta^2/L^2$. The factor-of-4 margin in
\eqref{eq:drift-eps} provides conservative headroom. $\square$
\end{proof}

For $\eta = 0.01$, $L \approx 1$: $\epsilon^* \approx 1.25 \times 10^{-5}$.
In the UNSW-NB15 simulation, adding 2,000 pseudo-labeled samples each round
produces KL drift consistently above this threshold (observed range:
$10^{-2}$ to $10^{+1}$), so the trigger does not skip retrains in this
setting. The benefit manifests in real deployment under stable traffic,
where the trigger avoids unnecessary retraining.

\subsection{EMA Normal Centroid}

AOC-IDS fixes the normal centroid to the initial training set,
preventing adaptation to benign traffic drift.

\textbf{Proposal.} At each online round $t$, update:
$\hat\mu_n^{(t+1)} = (1-\alpha_t)\hat\mu_n^{(t)} + \alpha_t\bar z_{\mathrm{pred\text{-}n}}^{(t)}$,
where $\bar z_{\mathrm{pred\text{-}n}}^{(t)}$ is the mean embedding of
samples predicted as normal in the current batch, renormalised to the unit sphere.

\begin{theorem}[EMA drift-tracking bound]\label{thm:ema}
If the true normal mean drifts at speed $\|\mu_n(t+1) - \mu_n(t)\| \le v$
and per-batch noise is $\sigma_z/\sqrt{B}$, then
$\alpha^* = \min\{1, \sqrt{vB/\sigma_z}\}$ gives
\begin{equation}
\mathbb{E}\|\hat\mu_n^{(t)} - \mu_n(t)\|^2 \le \frac{2\sigma_z v}{\sqrt{B}} + O(v^2),
\label{eq:ema-bound}
\end{equation}
which is order-optimal among constant-step EMA estimators.
\end{theorem}

\begin{proof}[Proof sketch]
MSE under step $\alpha$ decomposes as $\alpha^2\sigma_z^2/B$ (noise) $+
v^2/(2\alpha^2)$ (lag). Minimising over $\alpha$ gives $\alpha^* =
(vB/\sigma_z)^{1/2}$ and the stated bound. $\square$
\end{proof}

%%% ------------------------------------------------------------------ %%%
\section{Experiments}
%%% ------------------------------------------------------------------ %%%

\subsection{Setup}

We evaluate on UNSW-NB15~\cite{moustafa2015unsw} (196 features, MinMax
normalised). All hyperparameters other than the loss, architecture, and
pipeline match AOC-IDS exactly to isolate the contribution.
The run at 100 epochs is one-eighth of the paper's standard 800.

\medskip
\begin{table}[t]
\caption{Experimental Configuration}
\label{tab:config}
\centering
\begin{tabular}{lll}
\toprule
Parameter & CRC (Baseline) & GS-CRC (Ours)\\
\midrule
Loss & CRC & GS-CRC (Def.~\ref{def:gscrc})\\
Bottleneck $m$ & 64 & 32 (Prop.~\ref{prop:bottleneck})\\
Decision rule & Majority vote & LLR (Prop.~\ref{prop:llr})\\
Abnormal discriminator & 2-Gaussian & $K$-GMM, BIC $K \le 10$\\
Normal centroid & Fixed & EMA $\alpha=0.05$ (Thm.~\ref{thm:ema})\\
Update trigger & Fixed interval & KL-drift $\epsilon^*=10^{-5}$\\
Loss weighting & Equal & Learned $\sigma_e, \sigma_d$\\
Batch size & 256 & 1024\\
Epochs & 100 & 100\\
$\gamma$ & -- & 2.0\\
$\alpha_{\text{symm}}$ & -- & 0.5\\
$\beta_{\text{temp}}$ & -- & 0.5\\
$\mathcal{L}_{\max}$ & -- & 5.0\\
Optimizer & SGD, $10^{-3}$ & SGD, $10^{-3}$\\
PyTorch & 1.13.1 & 2.11.0+cu130\\
GPU & RTX 5060 Laptop & RTX 5060 Laptop\\
\bottomrule
\end{tabular}
\end{table}

Test set: 37,000 normal (44.9\%) and 45,332 attack (55.1\%).
Trivial all-attack baseline: accuracy 55.1\%, F1 0.710.
Trivial all-normal baseline: accuracy 44.9\%, F1 0.000.

\subsection{Per-Seed Results}

\begin{table}[t]
\caption{Per-Seed Results on UNSW-NB15 (100 epochs)}
\label{tab:per-seed}
\centering
\begin{tabular}{clcccc}
\toprule
& Seed & Acc & Prec & Rec & F1\\
\midrule
\multirow{7}{*}{\rotatebox{90}{GS-CRC (ours)}}
& 5009 & 0.8456 & 0.8092 & 0.9417 & 0.8704\\
& 5010 & 0.8409 & 0.7972 & 0.9536 & 0.8684\\
& 5011 & 0.8379 & 0.7852 & 0.9714 & 0.8684\\
& 5012 & 0.8320 & 0.7878 & 0.9510 & 0.8617\\
& 5013 & 0.8430 & 0.8015 & 0.9502 & 0.8695\\
\cmidrule{2-6}
& \textbf{Mean} & \textbf{0.8399} & \textbf{0.7962} & \textbf{0.9536} & \textbf{0.8677}\\
& Std & 0.0049 & 0.0087 & 0.0107 & 0.0032\\
\midrule
\multirow{7}{*}{\rotatebox{90}{CRC (baseline)}}
& 5009 & 0.5506 & 0.5506 & 1.0000 & 0.7102$^\dagger$\\
& 5010 & 0.4494 & 0.0000 & 0.0000 & 0.0000$^*$\\
& 5011 & 0.4827 & 1.0000 & 0.0605 & 0.1141$^*$\\
& 5012 & 0.5697 & 0.5749 & 0.8387 & 0.6822\\
& 5013 & 0.4494 & 0.0000 & 0.0000 & 0.0000$^*$\\
\cmidrule{2-6}
& Mean & 0.5004 & 0.4251 & 0.3798 & 0.3013\\
& Std & 0.0537 & 0.3905 & 0.4507 & 0.3151\\
\bottomrule
\end{tabular}
\smallskip\\
{\footnotesize $^\dagger$ Degenerate: all-attack. $^*$ Collapsed: all-normal.}
\end{table}

\subsection{Comparison Summary}

\begin{table}[t]
\caption{Summary Comparison (mean $\pm$ std, 5 seeds, 100 epochs)}
\label{tab:summary}
\centering
\begin{tabular}{lcccc}
\toprule
Method & Accuracy & Precision & Recall & F1\\
\midrule
CRC (base) & 0.500$\pm$.054 & 0.425$\pm$.391 & 0.380$\pm$.451 & 0.301$\pm$.315\\
GS-CRC & \textbf{0.840}$\pm$.005 & \textbf{0.796}$\pm$.009 & \textbf{0.954}$\pm$.011 & \textbf{0.868}$\pm$.003\\
\midrule
$\Delta$ (abs.) & +0.340 & +0.371 & +0.574 & +0.567\\
$\Delta$ (\%) & +68.0\% & +87.3\% & +151.1\% & +188.2\%\\
\bottomrule
\end{tabular}
\end{table}

\subsection{Confusion Matrices}

\begin{table}[t]
\caption{Confusion Matrices, Selected Seeds}
\label{tab:cm}
\centering
\begin{tabular}{lcccc}
\toprule
& \multicolumn{2}{c}{Pred Normal} & \multicolumn{2}{c}{Pred Attack}\\
Method & TN & FP & FN & TP\\
\midrule
CRC seed 5009 & 0 & 37000 & 0 & 45332\\
GS-CRC seed 5009 & 26933 & 10067 & 2644 & 42688\\
\midrule
CRC seed 5010 & 37000 & 0 & 45332 & 0\\
GS-CRC seed 5010 & 26003 & 10997 & 2103 & 43229\\
\midrule
CRC seed 5012 & 8884 & 28116 & 7312 & 38020\\
GS-CRC seed 5012 & 25391 & 11609 & 2223 & 43109\\
\bottomrule
\end{tabular}
\end{table}

\subsection{Analysis}

\noindent\textbf{Stability.}
All five GS-CRC seeds produce non-degenerate confusion matrices
(0/5 collapses, down from 3/5). Theorems~\ref{thm:fisher-crc}
and~\ref{thm:fisher-gscrc} predict exactly this: Fisher inconsistency
allows CRC to settle into degenerate geometry with positive probability,
while GS-CRC's $\mathcal{L}_a^+$ term rules it out.
F1 std drops from 0.315 to 0.003 --- a 99-fold stability improvement.

\noindent\textbf{Recall improvement.}
Mean recall rises from 0.380 to 0.954 (+151.1\%). The $K$-component GMM
correctly models the multi-modal abnormal score distribution (10 attack
categories produce 10 distinct cosine-similarity clusters), recovering the
$\Omega(1)$ excess risk identified in Theorem~\ref{thm:gmm-risk}. The
GS-CRC loss additionally ensures that minority-subclass representations
receive polynomial rather than exponential gradient, so all 10 attack
subtypes remain separated in the embedding space by round 100.

\noindent\textbf{Precision--recall tradeoff.}
Precision settles at $0.796 \pm 0.009$, down from $0.998 \pm 0.000$ in the
2-Gaussian version. The $K$-GMM models the abnormal distribution with more
components, raising $\log p(s|a)$ in overlap regions between low-frequency
attack subtypes and the normal class.
This shifts the LLR threshold, increasing sensitivity (recall) at the cost
of specificity (false positive rate $\approx 28\%$).
The false positive rate is controllable without retraining: increasing
$\log(\pi_n/\pi_a)$ in \eqref{eq:llr} shifts the operating point toward
higher precision along the precision--recall curve.

\noindent\textbf{Baseline artefact.}
CRC seed 5009 achieves F1 $= 0.710$ by predicting every sample as attack,
matching the trivial all-attack classifier exactly (recall $= 1$,
precision $= 0.551$). This is not a detection result.

%%% ------------------------------------------------------------------ %%%
\section{Limitations and Future Work}
%%% ------------------------------------------------------------------ %%%

\noindent\textbf{Epoch budget.}
All results are at 100 epochs (one-eighth of the standard 800).
The 800-epoch run on both CRC and GS-CRC is the required next experiment
for definitive numbers.

\noindent\textbf{Dataset scope.}
We evaluate only on UNSW-NB15. NSL-KDD evaluation is needed to confirm
that gradient-balance improvement transfers to $K=5$ attack types and
a different imbalance structure ($l_a^{(\min)}/l_a \approx 0.001$ for U2R).

\noindent\textbf{Parameter sweeps.}
$\gamma$, $\alpha_{\text{symm}}$, $\beta_{\text{temp}}$, and $\mathcal{L}_{\max}$
are set by theoretical argument rather than grid search. A full ablation
is left to the extended version.

\noindent\textbf{KL-drift trigger calibration.}
In the UNSW-NB15 simulation, adding 2,000 samples per round produces
KL drift consistently above $\epsilon^* = 10^{-5}$, so the trigger
retrains every round. For deployment under stable traffic the trigger
would skip unnecessary retrains; calibrating $\eta$ and $L$ for a specific
network environment is future work.

%%% ------------------------------------------------------------------ %%%
\section{Conclusion}
%%% ------------------------------------------------------------------ %%%

We presented a comprehensive upgrade to AOC-IDS addressing seven
structural pathologies. The GS-CRC loss fixes five defects in the CRC
objective via symmetric structure, focal modulation, adaptive temperature,
class-balance prefactors, and bounded influence. The detection pipeline
is upgraded with a $K$-component GMM discriminator (fixing $\Omega(1)$
excess risk), an EMA normal centroid (fixing static drift), a
KL-drift update trigger (fixing arbitrary cadence), and a calibrated
LLR decision rule.

On UNSW-NB15 at 100 epochs, the full system raises mean F1 from 0.301
to 0.868 (+188.2\%), recall from 0.380 to 0.954, eliminates all five
degenerate seeds, and reduces F1 standard deviation from 0.315 to 0.003.
Each improvement is backed by a formal theorem rather than an empirical
observation, making the contribution robust to dataset and deployment variation.

%%% ------------------------------------------------------------------ %%%
\section*{Acknowledgment}
We thank the authors of AOC-IDS for releasing their preprocessed datasets.

%%% ------------------------------------------------------------------ %%%
\begin{thebibliography}{20}

\bibitem{zhang2024aoc}
X.~Zhang, R.~Zhao, Z.~Jiang, Z.~Sun, Y.~Ding, E.~C.~H.~Ngai, and
S.-H.~Yang, ``AOC-IDS: Autonomous online framework with contrastive learning
for intrusion detection,'' in \textit{IEEE INFOCOM 2024}, pp.~581--590.

\bibitem{lin2017focal}
T.-Y.~Lin, P.~Goyal, R.~Girshick, K.~He, and P.~Dollár,
``Focal loss for dense object detection,'' in \textit{ICCV}, 2017,
pp.~2980--2988.

\bibitem{mirsky2018kitsune}
Y.~Mirsky, T.~Doitshman, Y.~Elovici, and A.~Shabtai,
``Kitsune: An ensemble of autoencoders for online network intrusion detection,''
in \textit{NDSS}, 2018.

\bibitem{ring2019flow}
M.~Ring, S.~Wunderlich, D.~Scheuring, D.~Landes, and A.~Hotho,
``A survey of network-based intrusion detection data sets,''
\textit{Computers \& Security}, vol.~86, pp.~147--167, 2019.

\bibitem{wang2021understanding}
T.~Wang and P.~Isola,
``Understanding contrastive representation learning through alignment and
uniformity on the hypersphere,'' in \textit{ICML}, 2021.

\bibitem{bartlett2007adaboost}
P.~L.~Bartlett and M.~H.~Tewari,
``Adaboost is consistent,''
\textit{J. Mach. Learn. Res.}, vol.~8, pp.~2347--2368, 2007.

\bibitem{huber1964robust}
P.~J.~Huber,
``Robust estimation of a location parameter,''
\textit{Ann. Math. Statist.}, vol.~35, no.~1, pp.~73--101, 1964.

\bibitem{natarajan2013learning}
N.~Natarajan, I.~S.~Dhillon, P.~K.~Ravikumar, and A.~Tewari,
``Learning with noisy labels,'' in \textit{NeurIPS}, 2013.

\bibitem{kendall2018multi}
A.~Kendall, Y.~Gal, and R.~Cipolla,
``Multi-task learning using uncertainty to weigh losses for scene geometry
and semantics,'' in \textit{CVPR}, 2018.

\bibitem{tishby2000information}
N.~Tishby, F.~C.~Pereira, and W.~Bialek,
``The information bottleneck method,'' in \textit{37th Allerton Conference},
2000.

\bibitem{moustafa2015unsw}
N.~Moustafa and J.~Slay,
``UNSW-NB15: A comprehensive data set for network intrusion detection systems,''
in \textit{MilCIS}, 2015.

\end{thebibliography}

\end{document}
